% WORKSHOP PAPER. 5 pages excluding references.
% Filename matches the Overleaf file so an upload overwrites it cleanly. Do not rename.
% The pristine template body is kept beside this as template_original.tex.
%
% Option 6, workshop with double-blind review. Verified in neurips_2026.sty:
% dblblindworkshop leaves \if@anonymous true, so the author block renders as
% "Anonymous Author(s)" automatically. sglblindworkshop would NOT.

\documentclass{article}

\usepackage[dblblindworkshop]{neurips_2026}

\usepackage[utf8]{inputenc}
\usepackage[T1]{fontenc}
\usepackage{hyperref}
\usepackage{url}
\usepackage{booktabs}
\usepackage{amsfonts}
\usepackage{nicefrac}
\usepackage{microtype}
\usepackage{xcolor}

% Every number lives in numbers.tex so this paper and long.tex cannot disagree.
% numbers.tex
%
% Every number either paper cites, defined once.
% Both neurips_2026.tex (workshop, 5 pages) and long.tex \input this file,
% so the two papers cannot drift apart.
%
% Rule: a macro is added here only after the figure exists in docs/paper_numbers.md
% with an n and a reproducing script. See that file's rule 1.
% The section number in each comment points into docs/paper_numbers.md.
%
% Naming: LaTeX control sequences cannot contain digits, so 3B is ThreeB
% and 7B is SevenB throughout.


% ---------------------------------------------------------------------------
% Splits and sample sizes
% ---------------------------------------------------------------------------

\newcommand{\nDev}{700}                  % testmini, development split
\newcommand{\nTest}{1{,}700}             % test.json, held-out split
\newcommand{\nIE}{600}
\newcommand{\nMATH}{600}
\newcommand{\nKNOW}{500}


% ---------------------------------------------------------------------------
% Headline accuracy, strict scoring, n = 1,700 held out. Section 2.18
% ---------------------------------------------------------------------------

\newcommand{\accThreeB}{59.8\%}
\newcommand{\accThreeBcount}{1016/1700}
\newcommand{\accThreeBci}{57.4--62.1}

\newcommand{\accSevenB}{67.8\%}
\newcommand{\accSevenBcount}{1153/1700}
\newcommand{\accSevenBci}{65.6--70.0}

\newcommand{\accRouted}{75.8\%}
\newcommand{\accRoutedcount}{1289/1700}
\newcommand{\accRoutedci}{73.7--77.8}

\newcommand{\accCloud}{77.4\%}
\newcommand{\accCloudcount}{1315/1700}
\newcommand{\accCloudci}{75.3--79.3}

\newcommand{\gainOverSevenB}{8.0}        % routed minus 7B alone
\newcommand{\gainOverThreeB}{16.0}       % routed minus 3B alone


% ---------------------------------------------------------------------------
% Routed against cloud alone, paired McNemar on identical claims. Section 2.18
% ---------------------------------------------------------------------------

\newcommand{\mcnemarAllP}{0.119}
\newcommand{\mcnemarAllDisagree}{258}
\newcommand{\mcnemarAllRoutedWins}{116}
\newcommand{\mcnemarAllCloudWins}{142}

\newcommand{\mcnemarIEP}{0.007}          % the one significant result, routed loses
\newcommand{\mcnemarIEDisagree}{109}
\newcommand{\mcnemarMATHP}{1.000}
\newcommand{\mcnemarKNOWP}{0.917}


% ---------------------------------------------------------------------------
% Per subset, all four arms. Section 2.18
% ---------------------------------------------------------------------------

\newcommand{\accThreeBIE}{63.3\%}
\newcommand{\accThreeBMATH}{57.3\%}
\newcommand{\accThreeBKNOW}{58.4\%}

\newcommand{\accSevenBIE}{68.5\%}
\newcommand{\accSevenBMATH}{66.5\%}
\newcommand{\accSevenBKNOW}{68.6\%}

\newcommand{\accRoutedIE}{77.5\%}
\newcommand{\accRoutedMATH}{77.3\%}
\newcommand{\accRoutedKNOW}{72.0\%}

\newcommand{\accCloudIE}{82.3\%}
\newcommand{\accCloudMATH}{77.2\%}
\newcommand{\accCloudKNOW}{71.6\%}

\newcommand{\escRateIE}{36.5\%}
\newcommand{\escRateMATH}{79.8\%}
\newcommand{\escRateKNOW}{42.0\%}


% ---------------------------------------------------------------------------
% Contribution 2: same-family agreement is correlated error. Section 2.18
% The whole FDV-IE deficit sits on claims the gate kept local.
% ---------------------------------------------------------------------------

\newcommand{\keptLocalIEn}{381}
\newcommand{\keptLocalIERouted}{76.9\%}
\newcommand{\keptLocalIECloud}{85.6\%}
\newcommand{\keptLocalIEGap}{8.7}        % points, routed below cloud

\newcommand{\escalatedIEn}{219}
\newcommand{\escalatedIERouted}{78.5\%}
\newcommand{\escalatedIECloud}{76.7\%}   % routed is slightly ahead here


% ---------------------------------------------------------------------------
% Escalation earns its cost. Section 2.18
% ---------------------------------------------------------------------------

\newcommand{\cloudCalls}{908}
\newcommand{\escWouldThreeB}{45.7\%}
\newcommand{\escWouldSevenB}{60.8\%}
\newcommand{\escCloudGave}{75.8\%}
\newcommand{\escGain}{15.0}              % points over the 7B on the same claims

\newcommand{\trigDisagreeN}{459}
\newcommand{\trigDisagreeSevenB}{57.3\%}
\newcommand{\trigDisagreeCloud}{73.4\%}
\newcommand{\trigDisagreeGain}{16.1}

\newcommand{\trigNumericN}{449}
\newcommand{\trigNumericSevenB}{64.4\%}
\newcommand{\trigNumericCloud}{78.2\%}
\newcommand{\trigNumericGain}{13.8}


% ---------------------------------------------------------------------------
% Routing behaviour. Section 2.18
% ---------------------------------------------------------------------------

\newcommand{\escalationRate}{53.4\%}     % claims that called the cloud
\newcommand{\onDeviceRate}{46.6\%}       % claims answered entirely on device
\newcommand{\localsAgreeRate}{66.2\%}
\newcommand{\accKeptLocal}{75.9\%}
\newcommand{\accCloudCalled}{75.8\%}


% ---------------------------------------------------------------------------
% Cost, from our own recorded token counts. Section 2.18, and 2.18.1 and 3.5
% USD and tokens only. Never CNY.
% ---------------------------------------------------------------------------

\newcommand{\costRouted}{\$1.066}
\newcommand{\costCloud}{\$1.882}
\newcommand{\costRatio}{56.6\%}
\newcommand{\tokensRoutedIn}{3{,}400{,}854}
\newcommand{\tokensRoutedOut}{2{,}106{,}383}
\newcommand{\tokensCloudIn}{6{,}319{,}626}
\newcommand{\tokensCloudOut}{3{,}560{,}933}
\newcommand{\medianLatency}{26.0}        % seconds per claim, GPU box, report median not mean


% ---------------------------------------------------------------------------
% Oracles. Upper bounds only. Neither may be presented as a system result.
% Section 2.18
% ---------------------------------------------------------------------------

\newcommand{\oracleThreeWay}{84.8\%}
\newcommand{\oracleThreeWayci}{83.0--86.4}
\newcommand{\oracleLocalOnly}{79.9\%}    % beats the cloud with zero cloud calls
\newcommand{\oracleLocalOnlyci}{78.0--81.8}


% ---------------------------------------------------------------------------
% Unparseable output. Section 2.18
% Cloud denominator is calls made, not 1,700.
% ---------------------------------------------------------------------------

\newcommand{\unparseThreeB}{1.65\%}
\newcommand{\unparseSevenB}{2.82\%}
\newcommand{\unparseCloud}{1.65\%}
\newcommand{\unparseRouted}{0.94\%}


% ---------------------------------------------------------------------------
% Retrieval, macro recall at k = 10. Section 1 and 2.16
% These are retrieval results. They may NOT be written as accuracy drivers.
% ---------------------------------------------------------------------------

\newcommand{\recallOursDev}{74.60\%}
\newcommand{\recallOursTest}{75.16\%}
\newcommand{\recallEmbedDev}{68.01\%}    % text-embedding-3-large, best published
\newcommand{\recallEmbedTest}{69.75\%}
\newcommand{\recallUpstreamBMDev}{65.16\%}
\newcommand{\recallUpstreamBMTest}{64.29\%}
\newcommand{\recallContrieverDev}{33.48\%}
\newcommand{\recallContrieverTest}{35.02\%}

\newcommand{\recallGainDev}{6.6}         % points over best published, testmini
\newcommand{\recallGainTest}{5.4}        % points over best published, held out

\newcommand{\recallElementTest}{70.0\%}
\newcommand{\recallAllGoldTest}{49.4\%}

\newcommand{\recallOursIETest}{77.92\%}
\newcommand{\recallOursMATHTest}{80.92\%}
\newcommand{\recallOursKNOWTest}{64.94\%}

% Why ours beats upstream's BM25: an implementation gap, priced.
\newcommand{\ablationIDF}{2.51}          % Lucene IDF, points
\newcommand{\ablationPunct}{2.02}        % dropping punctuation tokens, points
\newcommand{\ablationStem}{0.06}         % stemming, worth nothing
\newcommand{\tableInflation}{1.81}       % punctuation inflates table length this much
\newcommand{\paraInflation}{1.13}        % against this for paragraphs

% Fusion is dead on both splits.
\newcommand{\recallFusionTest}{74.95\%}  % below our BM25 alone


% ---------------------------------------------------------------------------
% Contribution 4: the error taxonomy. Section 2.17
% ---------------------------------------------------------------------------

\newcommand{\missingInfoKnowDev}{63.5\%}
\newcommand{\missingInfoKnowTest}{63.4\%}
\newcommand{\missingInfoAllTest}{41.0\%}
\newcommand{\missingInfoCiteRateTest}{38.4\%}
\newcommand{\missingInfoWrongTest}{47.6\%}
\newcommand{\missingInfoBaseTest}{35.8\%}


% ---------------------------------------------------------------------------
% The refuted bias, replicated four times. Section 2.18
% Share of claims each arm answers "entailed", where gold is 50.1% entailed.
% ---------------------------------------------------------------------------

\newcommand{\saysEntailedThreeB}{39.5\%}
\newcommand{\saysEntailedSevenB}{49.2\%}  % the only nearly calibrated arm
\newcommand{\saysEntailedRouted}{41.1\%}
\newcommand{\saysEntailedCloud}{33.9\%}
\newcommand{\saysEntailedGold}{50.1\%}

\newcommand{\entailedRouted}{67.1\%}
\newcommand{\entailedCloud}{61.7\%}
\newcommand{\refutedRouted}{84.6\%}
\newcommand{\refutedCloud}{93.1\%}


% ---------------------------------------------------------------------------
% The target device. Hard constraint, not a result. CLAUDE.md and section 3.1
% ---------------------------------------------------------------------------

\newcommand{\deviceLatency}{7}           % minutes per example, 3B on the MacBook
\newcommand{\deviceSecPerToken}{0.0641}
\newcommand{\deviceRsq}{0.995}
\newcommand{\deviceMeanPrompt}{6{,}610}  % tokens


% TITLE IS A PLACEHOLDER AND NEEDS YOUR APPROVAL.
\title{Routing Financial Claim Verification Between On-Device and Cloud Models}

% Fills the first-page footnote: "... (NeurIPS 2026). Workshop: <this>."
% Was unset since 10 August.
\workshoptitle{On-Device Intelligence: Foundation Models under Real-World Constraints}

% Double-blind. The style file replaces this block with "Anonymous Author(s)",
% so nothing here reaches the PDF. It is filled in for the camera-ready only.
\author{%
  Anonymous Author(s) \\
  Affiliation \\
  Address \\
  \texttt{email} \\
}

\begin{document}

\maketitle

\begin{abstract}
Verifying a financial claim against a regulatory filing requires reading a document far
longer than a small model's context window, and sending every filing to a frontier model
is expensive. We build a claim verification pipeline that runs two small models on device
and calls a cloud model only when a disagreement gate or an arithmetic detector fires. On
\nTest{} held-out claims from FINDVER the routed pipeline scores \accRouted{} against a
cloud-only baseline's \accCloud{}. The difference is not significant (paired McNemar,
\mcnemarAllDisagree{} disagreements, $p = \mcnemarAllP{}$), and the routed system uses
\costRatio{} of the cloud tokens with \onDeviceRate{} of claims answered entirely on
device. That aggregate hides structure that matters to anyone deploying it. The tie is a
significant loss on one subset offset by ties on the others, and we identify the cause:
the gate reads agreement between two models of the same family as confidence, and on that
subset the agreement is correlated error worth \keptLocalIEGap{} points. We report a BM25
retriever that reaches \recallOursTest{} macro recall with no neural model and no API
call, \recallGainTest{} points above the strongest retriever published on this benchmark,
and an error taxonomy in which one failure mode accounts for \missingInfoKnowTest{} of
knowledge-subset errors.
\end{abstract}

\section{Introduction}
\label{sec:intro}

% TODO. ~0.75 pages. Problem, setting, and the four contributions in this order:
%   1. routed pipeline ties frontier cloud at 57% of the cost, no training
%   2. same-family model agreement is correlated error
%   3. unconditional escalation of arithmetic claims beats a disagreement gate, and why
%   4. an error taxonomy on the held-out split
% Framing settled 21 Aug: move 1 is the system result, move 2 is the check that makes it
% credible. See working_state.md.

\section{Related work}
\label{sec:related}

% TODO. ~0.5 pages. FINDVER, MACE, RAG, edge-cloud routing.
% Phrasing constraint: "we found no other method evaluated on FINDVER." Never "nobody has."

\section{Method}
\label{sec:method}

% TODO. ~1.1 pages. BM25 retriever, the two local models, the gate, the numeric detector.

% Retrieval. Numbers from numbers.tex; provenance in docs/paper_numbers.md
% sections 1, 1.2, 1.5, 1.7, 2.8, 2.8.2, 2.16 and 4.5.
\subsection{Retrieval}
\label{sec:retrieval}

Reports in this benchmark are too long to fit in a prompt, so every system evaluated on
it retrieves first. Our retriever is BM25 over the report's own context elements. It runs
no neural model and makes no API call. On the held-out split it reaches \recallOursTest{}
macro recall at $k{=}10$, against \recallEmbedTest{} for \texttt{text-embedding-3-large},
the strongest retriever reported on this benchmark, and \recallUpstreamBMTest{} for the
benchmark's own BM25. Element recall is \recallElementTest{} and \recallAllGoldTest{} of
claims receive every element they need.

The eleven point gap to the benchmark's own BM25 is implementation rather than
algorithm, and we can account for it on the development split. Lucene's IDF variant is
worth \ablationIDF{} points and discarding punctuation tokens \ablationPunct{}. Stemming
is worth nothing. The punctuation result has a mechanism behind it: treating punctuation
as tokens inflates measured element length by \tableInflation{}$\times$ for tables against
\paraInflation{}$\times$ for paragraphs, so BM25's length normalisation penalises tables
and pushes them below the cut. Tables are roughly 18\% of elements and carry much of this
benchmark's evidence. We report this as a measurement, not as a design insight. Our
tokenizer discards punctuation because it was reused from an evidence-checking utility,
and the mechanism was found afterwards.

Three upgrades we expected to help did not. Decomposing a claim into sub-claims and merging
six ways of ranking never beat querying with the whole claim, and at matched candidate
count the single query wins by 3.2 points. Fusing a dense arm into the ranking gains
nothing untuned, and a weak arm actively harms: reciprocal rank fusion gives every list an
equal vote, so adding a 33\% retriever costs 5.1 points. Sweeping $k_1$ and $b$ moves the
score inside its own noise.

We report these as retrieval results only. In our pipeline, retrieval quality and answer
quality are close to decoupled. Supplying perfect evidence does not raise accuracy
($p = 0.116$), removing a third of the evidence does not lower it ($p = 1.000$), and
raising $k$ from 10 to 20 buys 10.8 points of recall and no measurable accuracy at 87\%
more prompt tokens. A recall gain on this benchmark should not be assumed to be a
verification gain.

\section{Experimental setup}
\label{sec:setup}

% TODO. ~0.4 pages. Splits, strict scoring with the unparseable rate beside it,
% and the testmini/test protocol: everything tuned on testmini, test.json touched once.

\section{Results}
\label{sec:results}

% TODO. ~1.6 pages. The one surviving table goes here (four arms by subset), then
% escalation value, the FDV-IE mechanism, and cost as prose.
% The 16 published baselines and the retriever comparison table are cut from this paper
% and live in long.tex.

\section{Limitations and future work}
\label{sec:limitations}

% TODO. ~0.4 pages. Tuning disclosure, the FDV-IE gate leak as a finding never fixed,
% and the cloud model ignoring seed and temperature.

% \section*{NeurIPS Paper Checklist}

%%% BEGIN INSTRUCTIONS %%%
The checklist is designed to encourage best practices for responsible machine learning research, addressing issues of reproducibility, transparency, research ethics, and societal impact. Do not remove the checklist: {\bf The papers not including the checklist will be desk rejected.} The checklist should follow the references and follow the (optional) supplemental material.  The checklist does NOT count towards the page
limit. 

Please read the checklist guidelines carefully for information on how to answer these questions. For each question in the checklist:
\begin{itemize}
    \item You should answer \answerYes{}, \answerNo{}, or \answerNA{}.
    \item \answerNA{} means either that the question is Not Applicable for that particular paper or the relevant information is Not Available.
    \item Please provide a short (1--2 sentence) justification right after your answer (even for \answerNA). 
   % \item {\bf The papers not including the checklist will be desk rejected.}
\end{itemize}

{\bf The checklist answers are an integral part of your paper submission.} They are visible to the reviewers, area chairs, senior area chairs, and ethics reviewers. You will also be asked to include it (after eventual revisions) with the final version of your paper, and its final version will be published with the paper.

The reviewers of your paper will be asked to use the checklist as one of the factors in their evaluation. While \answerYes{} is generally preferable to \answerNo{}, it is perfectly acceptable to answer \answerNo{} provided a proper justification is given (e.g., error bars are not reported because it would be too computationally expensive'' or ``we were unable to find the license for the dataset we used''). In general, answering \answerNo{} or \answerNA{} is not grounds for rejection. While the questions are phrased in a binary way, we acknowledge that the true answer is often more nuanced, so please just use your best judgment and write a justification to elaborate. All supporting evidence can appear either in the main paper or the supplemental material, provided in appendix. If you answer \answerYes{} to a question, in the justification please point to the section(s) where related material for the question can be found.

IMPORTANT, please:
\begin{itemize}
    \item {\bf Delete this instruction block, but keep the section heading ``NeurIPS Paper Checklist"},
    \item  {\bf Keep the checklist subsection headings, questions/answers and guidelines below.}
    \item {\bf Do not modify the questions and only use the provided macros for your answers}.
\end{itemize} 
 

%%% END INSTRUCTIONS %%%


\begin{enumerate}

\item {\bf Claims}
    \item[] Question: Do the main claims made in the abstract and introduction accurately reflect the paper's contributions and scope?
    \item[] Answer: \answerTODO{} % Replace by \answerYes{}, \answerNo{}, or \answerNA{}.
    \item[] Justification: \justificationTODO{}
    \item[] Guidelines:
    \begin{itemize}
        \item The answer \answerNA{} means that the abstract and introduction do not include the claims made in the paper.
        \item The abstract and/or introduction should clearly state the claims made, including the contributions made in the paper and important assumptions and limitations. A \answerNo{} or \answerNA{} answer to this question will not be perceived well by the reviewers. 
        \item The claims made should match theoretical and experimental results, and reflect how much the results can be expected to generalize to other settings. 
        \item It is fine to include aspirational goals as motivation as long as it is clear that these goals are not attained by the paper. 
    \end{itemize}

\item {\bf Limitations}
    \item[] Question: Does the paper discuss the limitations of the work performed by the authors?
    \item[] Answer: \answerTODO{} % Replace by \answerYes{}, \answerNo{}, or \answerNA{}.
    \item[] Justification: \justificationTODO{}
    \item[] Guidelines:
    \begin{itemize}
        \item The answer \answerNA{} means that the paper has no limitation while the answer \answerNo{} means that the paper has limitations, but those are not discussed in the paper. 
        \item The authors are encouraged to create a separate ``Limitations'' section in their paper.
        \item The paper should point out any strong assumptions and how robust the results are to violations of these assumptions (e.g., independence assumptions, noiseless settings, model well-specification, asymptotic approximations only holding locally). The authors should reflect on how these assumptions might be violated in practice and what the implications would be.
        \item The authors should reflect on the scope of the claims made, e.g., if the approach was only tested on a few datasets or with a few runs. In general, empirical results often depend on implicit assumptions, which should be articulated.
        \item The authors should reflect on the factors that influence the performance of the approach. For example, a facial recognition algorithm may perform poorly when image resolution is low or images are taken in low lighting. Or a speech-to-text system might not be used reliably to provide closed captions for online lectures because it fails to handle technical jargon.
        \item The authors should discuss the computational efficiency of the proposed algorithms and how they scale with dataset size.
        \item If applicable, the authors should discuss possible limitations of their approach to address problems of privacy and fairness.
        \item While the authors might fear that complete honesty about limitations might be used by reviewers as grounds for rejection, a worse outcome might be that reviewers discover limitations that aren't acknowledged in the paper. The authors should use their best judgment and recognize that individual actions in favor of transparency play an important role in developing norms that preserve the integrity of the community. Reviewers will be specifically instructed to not penalize honesty concerning limitations.
    \end{itemize}

\item {\bf Theory assumptions and proofs}
    \item[] Question: For each theoretical result, does the paper provide the full set of assumptions and a complete (and correct) proof?
    \item[] Answer: \answerTODO{} % Replace by \answerYes{}, \answerNo{}, or \answerNA{}.
    \item[] Justification: \justificationTODO{}
    \item[] Guidelines:
    \begin{itemize}
        \item The answer \answerNA{} means that the paper does not include theoretical results. 
        \item All the theorems, formulas, and proofs in the paper should be numbered and cross-referenced.
        \item All assumptions should be clearly stated or referenced in the statement of any theorems.
        \item The proofs can either appear in the main paper or the supplemental material, but if they appear in the supplemental material, the authors are encouraged to provide a short proof sketch to provide intuition. 
        \item Inversely, any informal proof provided in the core of the paper should be complemented by formal proofs provided in appendix or supplemental material.
        \item Theorems and Lemmas that the proof relies upon should be properly referenced. 
    \end{itemize}

    \item {\bf Experimental result reproducibility}
    \item[] Question: Does the paper fully disclose all the information needed to reproduce the main experimental results of the paper to the extent that it affects the main claims and/or conclusions of the paper (regardless of whether the code and data are provided or not)?
    \item[] Answer: \answerTODO{} % Replace by \answerYes{}, \answerNo{}, or \answerNA{}.
    \item[] Justification: \justificationTODO{}
    \item[] Guidelines:
    \begin{itemize}
        \item The answer \answerNA{} means that the paper does not include experiments.
        \item If the paper includes experiments, a \answerNo{} answer to this question will not be perceived well by the reviewers: Making the paper reproducible is important, regardless of whether the code and data are provided or not.
        \item If the contribution is a dataset and\slash or model, the authors should describe the steps taken to make their results reproducible or verifiable. 
        \item Depending on the contribution, reproducibility can be accomplished in various ways. For example, if the contribution is a novel architecture, describing the architecture fully might suffice, or if the contribution is a specific model and empirical evaluation, it may be necessary to either make it possible for others to replicate the model with the same dataset, or provide access to the model. In general. releasing code and data is often one good way to accomplish this, but reproducibility can also be provided via detailed instructions for how to replicate the results, access to a hosted model (e.g., in the case of a large language model), releasing of a model checkpoint, or other means that are appropriate to the research performed.
        \item While NeurIPS does not require releasing code, the conference does require all submissions to provide some reasonable avenue for reproducibility, which may depend on the nature of the contribution. For example
        \begin{enumerate}
            \item If the contribution is primarily a new algorithm, the paper should make it clear how to reproduce that algorithm.
            \item If the contribution is primarily a new model architecture, the paper should describe the architecture clearly and fully.
            \item If the contribution is a new model (e.g., a large language model), then there should either be a way to access this model for reproducing the results or a way to reproduce the model (e.g., with an open-source dataset or instructions for how to construct the dataset).
            \item We recognize that reproducibility may be tricky in some cases, in which case authors are welcome to describe the particular way they provide for reproducibility. In the case of closed-source models, it may be that access to the model is limited in some way (e.g., to registered users), but it should be possible for other researchers to have some path to reproducing or verifying the results.
        \end{enumerate}
    \end{itemize}


\item {\bf Open access to data and code}
    \item[] Question: Does the paper provide open access to the data and code, with sufficient instructions to faithfully reproduce the main experimental results, as described in supplemental material?
    \item[] Answer: \answerTODO{} % Replace by \answerYes{}, \answerNo{}, or \answerNA{}.
    \item[] Justification: \justificationTODO{}
    \item[] Guidelines:
    \begin{itemize}
        \item The answer \answerNA{} means that paper does not include experiments requiring code.
        \item Please see the NeurIPS code and data submission guidelines (\url{https://neurips.cc/public/guides/CodeSubmissionPolicy}) for more details.
        \item While we encourage the release of code and data, we understand that this might not be possible, so \answerNo{} is an acceptable answer. Papers cannot be rejected simply for not including code, unless this is central to the contribution (e.g., for a new open-source benchmark).
        \item The instructions should contain the exact command and environment needed to run to reproduce the results. See the NeurIPS code and data submission guidelines (\url{https://neurips.cc/public/guides/CodeSubmissionPolicy}) for more details.
        \item The authors should provide instructions on data access and preparation, including how to access the raw data, preprocessed data, intermediate data, and generated data, etc.
        \item The authors should provide scripts to reproduce all experimental results for the new proposed method and baselines. If only a subset of experiments are reproducible, they should state which ones are omitted from the script and why.
        \item At submission time, to preserve anonymity, the authors should release anonymized versions (if applicable).
        \item Providing as much information as possible in supplemental material (appended to the paper) is recommended, but including URLs to data and code is permitted.
    \end{itemize}


\item {\bf Experimental setting/details}
    \item[] Question: Does the paper specify all the training and test details (e.g., data splits, hyperparameters, how they were chosen, type of optimizer) necessary to understand the results?
    \item[] Answer: \answerTODO{} % Replace by \answerYes{}, \answerNo{}, or \answerNA{}.
    \item[] Justification: \justificationTODO{}
    \item[] Guidelines:
    \begin{itemize}
        \item The answer \answerNA{} means that the paper does not include experiments.
        \item The experimental setting should be presented in the core of the paper to a level of detail that is necessary to appreciate the results and make sense of them.
        \item The full details can be provided either with the code, in appendix, or as supplemental material.
    \end{itemize}

\item {\bf Experiment statistical significance}
    \item[] Question: Does the paper report error bars suitably and correctly defined or other appropriate information about the statistical significance of the experiments?
    \item[] Answer: \answerTODO{} % Replace by \answerYes{}, \answerNo{}, or \answerNA{}.
    \item[] Justification: \justificationTODO{}
    \item[] Guidelines:
    \begin{itemize}
        \item The answer \answerNA{} means that the paper does not include experiments.
        \item The authors should answer \answerYes{} if the results are accompanied by error bars, confidence intervals, or statistical significance tests, at least for the experiments that support the main claims of the paper.
        \item The factors of variability that the error bars are capturing should be clearly stated (for example, train/test split, initialization, random drawing of some parameter, or overall run with given experimental conditions).
        \item The method for calculating the error bars should be explained (closed form formula, call to a library function, bootstrap, etc.)
        \item The assumptions made should be given (e.g., Normally distributed errors).
        \item It should be clear whether the error bar is the standard deviation or the standard error of the mean.
        \item It is OK to report 1-sigma error bars, but one should state it. The authors should preferably report a 2-sigma error bar than state that they have a 96\% CI, if the hypothesis of Normality of errors is not verified.
        \item For asymmetric distributions, the authors should be careful not to show in tables or figures symmetric error bars that would yield results that are out of range (e.g., negative error rates).
        \item If error bars are reported in tables or plots, the authors should explain in the text how they were calculated and reference the corresponding figures or tables in the text.
    \end{itemize}

\item {\bf Experiments compute resources}
    \item[] Question: For each experiment, does the paper provide sufficient information on the computer resources (type of compute workers, memory, time of execution) needed to reproduce the experiments?
    \item[] Answer: \answerTODO{} % Replace by \answerYes{}, \answerNo{}, or \answerNA{}.
    \item[] Justification: \justificationTODO{}
    \item[] Guidelines:
    \begin{itemize}
        \item The answer \answerNA{} means that the paper does not include experiments.
        \item The paper should indicate the type of compute workers CPU or GPU, internal cluster, or cloud provider, including relevant memory and storage.
        \item The paper should provide the amount of compute required for each of the individual experimental runs as well as estimate the total compute. 
        \item The paper should disclose whether the full research project required more compute than the experiments reported in the paper (e.g., preliminary or failed experiments that didn't make it into the paper). 
    \end{itemize}
    
\item {\bf Code of ethics}
    \item[] Question: Does the research conducted in the paper conform, in every respect, with the NeurIPS Code of Ethics \url{https://neurips.cc/public/EthicsGuidelines}?
    \item[] Answer: \answerTODO{} % Replace by \answerYes{}, \answerNo{}, or \answerNA{}.
    \item[] Justification: \justificationTODO{}
    \item[] Guidelines:
    \begin{itemize}
        \item The answer \answerNA{} means that the authors have not reviewed the NeurIPS Code of Ethics.
        \item If the authors answer \answerNo, they should explain the special circumstances that require a deviation from the Code of Ethics.
        \item The authors should make sure to preserve anonymity (e.g., if there is a special consideration due to laws or regulations in their jurisdiction).
    \end{itemize}


\item {\bf Broader impacts}
    \item[] Question: Does the paper discuss both potential positive societal impacts and negative societal impacts of the work performed?
    \item[] Answer: \answerTODO{} % Replace by \answerYes{}, \answerNo{}, or \answerNA{}.
    \item[] Justification: \justificationTODO{}
    \item[] Guidelines:
    \begin{itemize}
        \item The answer \answerNA{} means that there is no societal impact of the work performed.
        \item If the authors answer \answerNA{} or \answerNo, they should explain why their work has no societal impact or why the paper does not address societal impact.
        \item Examples of negative societal impacts include potential malicious or unintended uses (e.g., disinformation, generating fake profiles, surveillance), fairness considerations (e.g., deployment of technologies that could make decisions that unfairly impact specific groups), privacy considerations, and security considerations.
        \item The conference expects that many papers will be foundational research and not tied to particular applications, let alone deployments. However, if there is a direct path to any negative applications, the authors should point it out. For example, it is legitimate to point out that an improvement in the quality of generative models could be used to generate Deepfakes for disinformation. On the other hand, it is not needed to point out that a generic algorithm for optimizing neural networks could enable people to train models that generate Deepfakes faster.
        \item The authors should consider possible harms that could arise when the technology is being used as intended and functioning correctly, harms that could arise when the technology is being used as intended but gives incorrect results, and harms following from (intentional or unintentional) misuse of the technology.
        \item If there are negative societal impacts, the authors could also discuss possible mitigation strategies (e.g., gated release of models, providing defenses in addition to attacks, mechanisms for monitoring misuse, mechanisms to monitor how a system learns from feedback over time, improving the efficiency and accessibility of ML).
    \end{itemize}
    
\item {\bf Safeguards}
    \item[] Question: Does the paper describe safeguards that have been put in place for responsible release of data or models that have a high risk for misuse (e.g., pre-trained language models, image generators, or scraped datasets)?
    \item[] Answer: \answerTODO{} % Replace by \answerYes{}, \answerNo{}, or \answerNA{}.
    \item[] Justification: \justificationTODO{}
    \item[] Guidelines:
    \begin{itemize}
        \item The answer \answerNA{} means that the paper poses no such risks.
        \item Released models that have a high risk for misuse or dual-use should be released with necessary safeguards to allow for controlled use of the model, for example by requiring that users adhere to usage guidelines or restrictions to access the model or implementing safety filters. 
        \item Datasets that have been scraped from the Internet could pose safety risks. The authors should describe how they avoided releasing unsafe images.
        \item We recognize that providing effective safeguards is challenging, and many papers do not require this, but we encourage authors to take this into account and make a best faith effort.
    \end{itemize}

\item {\bf Licenses for existing assets}
    \item[] Question: Are the creators or original owners of assets (e.g., code, data, models), used in the paper, properly credited and are the license and terms of use explicitly mentioned and properly respected?
    \item[] Answer: \answerTODO{} % Replace by \answerYes{}, \answerNo{}, or \answerNA{}.
    \item[] Justification: \justificationTODO{}
    \item[] Guidelines:
    \begin{itemize}
        \item The answer \answerNA{} means that the paper does not use existing assets.
        \item The authors should cite the original paper that produced the code package or dataset.
        \item The authors should state which version of the asset is used and, if possible, include a URL.
        \item The name of the license (e.g., CC-BY 4.0) should be included for each asset.
        \item For scraped data from a particular source (e.g., website), the copyright and terms of service of that source should be provided.
        \item If assets are released, the license, copyright information, and terms of use in the package should be provided. For popular datasets, \url{paperswithcode.com/datasets} has curated licenses for some datasets. Their licensing guide can help determine the license of a dataset.
        \item For existing datasets that are re-packaged, both the original license and the license of the derived asset (if it has changed) should be provided.
        \item If this information is not available online, the authors are encouraged to reach out to the asset's creators.
    \end{itemize}

\item {\bf New assets}
    \item[] Question: Are new assets introduced in the paper well documented and is the documentation provided alongside the assets?
    \item[] Answer: \answerTODO{} % Replace by \answerYes{}, \answerNo{}, or \answerNA{}.
    \item[] Justification: \justificationTODO{}
    \item[] Guidelines:
    \begin{itemize}
        \item The answer \answerNA{} means that the paper does not release new assets.
        \item Researchers should communicate the details of the dataset\slash code\slash model as part of their submissions via structured templates. This includes details about training, license, limitations, etc. 
        \item The paper should discuss whether and how consent was obtained from people whose asset is used.
        \item At submission time, remember to anonymize your assets (if applicable). You can either create an anonymized URL or include an anonymized zip file.
    \end{itemize}

\item {\bf Crowdsourcing and research with human subjects}
    \item[] Question: For crowdsourcing experiments and research with human subjects, does the paper include the full text of instructions given to participants and screenshots, if applicable, as well as details about compensation (if any)? 
    \item[] Answer: \answerTODO{} % Replace by \answerYes{}, \answerNo{}, or \answerNA{}.
    \item[] Justification: \justificationTODO{}
    \item[] Guidelines:
    \begin{itemize}
        \item The answer \answerNA{} means that the paper does not involve crowdsourcing nor research with human subjects.
        \item Including this information in the supplemental material is fine, but if the main contribution of the paper involves human subjects, then as much detail as possible should be included in the main paper. 
        \item According to the NeurIPS Code of Ethics, workers involved in data collection, curation, or other labor should be paid at least the minimum wage in the country of the data collector. 
    \end{itemize}

\item {\bf Institutional review board (IRB) approvals or equivalent for research with human subjects}
    \item[] Question: Does the paper describe potential risks incurred by study participants, whether such risks were disclosed to the subjects, and whether Institutional Review Board (IRB) approvals (or an equivalent approval/review based on the requirements of your country or institution) were obtained?
    \item[] Answer: \answerTODO{} % Replace by \answerYes{}, \answerNo{}, or \answerNA{}.
    \item[] Justification: \justificationTODO{}
    \item[] Guidelines:
    \begin{itemize}
        \item The answer \answerNA{} means that the paper does not involve crowdsourcing nor research with human subjects.
        \item Depending on the country in which research is conducted, IRB approval (or equivalent) may be required for any human subjects research. If you obtained IRB approval, you should clearly state this in the paper. 
        \item We recognize that the procedures for this may vary significantly between institutions and locations, and we expect authors to adhere to the NeurIPS Code of Ethics and the guidelines for their institution. 
        \item For initial submissions, do not include any information that would break anonymity (if applicable), such as the institution conducting the review.
    \end{itemize}

\item {\bf Declaration of LLM usage}
    \item[] Question: Does the paper describe the usage of LLMs if it is an important, original, or non-standard component of the core methods in this research? Note that if the LLM is used only for writing, editing, or formatting purposes and does \emph{not} impact the core methodology, scientific rigor, or originality of the research, declaration is not required.
    %this research? 
    \item[] Answer: \answerTODO{} % Replace by \answerYes{}, \answerNo{}, or \answerNA{}.
    \item[] Justification: \justificationTODO{}
    \item[] Guidelines:
    \begin{itemize}
        \item The answer \answerNA{} means that the core method development in this research does not involve LLMs as any important, original, or non-standard components.
        \item Please refer to our LLM policy in the NeurIPS handbook for what should or should not be described.
    \end{itemize}

\end{enumerate}
% Commented out. checklist.tex is the NeurIPS main-conference questionnaire and the
% workshop probably does not require it. UNVERIFIED: check the workshop's call for papers
% before submitting. Template says checklist pages do not count toward the limit anyway.

\end{document}
